\section{Discussion and Conclusion}\label{sec:conclusion}

Across 90+ experiments and 6000+ evaluation episodes on two VLA model families, three findings converge on a single conclusion: current VLA safety evaluation is structurally unreliable. Action collapse renders reported CAR illusory (Pi0+A1 drops from 91.5\% to ${\approx}$\,0\% under displacement filtering); EE-only oracle refinement is counterproductive, \emph{reducing} CAR below unrefined baselines; and collision-free demonstrations, not geometric supervision, account for all measured avoidance gains in fine-tuned models. We discuss the mechanistic basis for each finding and its implications.

\paragraph{Why EE-only SDF is redundant: a two-level account.}
The flat $\lambda$-CAR curve ($\lambda \in [0.005, 0.1]$, 20$\times$ range, all within the null band) rules out insufficient loss weight as the explanation and points to two independent structural causes.

\textit{(i) Demonstration data subsumes the SDF training signal.} Collision-free demonstrations shift the action distribution away from collision-prone trajectories by construction: every demonstration trajectory is safe, so imitating demonstrations implicitly encodes spatial avoidance. The ${\sim}$28\,pp CAR gain from LoRA fine-tuning (26.8\% $\to$ 54.5\%) arises entirely from this distributional shift, with no SDF contribution. The per-task CAR ordering (chocolate\_pudding highest, bbq\_sauce lowest) is invariant across all SDF conditions and matches demonstration coverage, confirming that avoidance patterns originate from data statistics rather than geometric gradients. In this regime, EE-only SDF supplies a signal the model has already learned from imitation; adding a redundant gradient to a converged behavior produces no measurable effect. The Phase~2 ablation ($\lambda = 0$ after step~5K, CAR = 57.5\%) reinforces this interpretation: removing SDF supervision mid-training leaves avoidance unchanged, because the model never relied on it.

\textit{(ii) End-effector monitoring cannot prevent full-body collisions.} Even granting a perfect distance signal, the SDF gradient acts on a single point (the EE) while collisions predominantly occur along intermediate arm links and joints. In SafeLIBERO manipulation tasks, the robot must reach around obstacles, placing the forearm and wrist in close proximity to obstacle surfaces that the EE-only SDF does not monitor. Oracle experiments confirm this structural gap: ground-truth SDF refinement \emph{reduces} CAR by 6\textendash10\,pp below unrefined baselines (73.1\textendash76.9\% vs.\ 82\textendash83\%), because pushing the EE away from obstacles rotates unmonitored links through them. The action space compounds the mismatch: VLA policies output delta EE poses, so the SDF gradient operates in the same low-dimensional space as the action, with no mechanism to constrain the full joint configuration that realizes a given EE displacement.

These two levels are independent. Resolving (i) by finding a training signal beyond what demonstrations provide would not address (ii), and resolving (ii) through full-body monitoring would not change the fact that demonstrations already supply the avoidance signal. Effective geometric safety for VLA manipulation requires advances on both fronts: supervision that covers the full kinematic chain, and training signals that provide collision information beyond what is already implicit in collision-free demonstrations.

\paragraph{Implications for VLA safety evaluation.}
CAR$_\text{f}$ corrects action-collapse contamination using only end-effector position logging, already available in most evaluation pipelines. The bimodal displacement distribution (collapsed below 1\,$\mu$m, active above 200\,mm) makes $\tau_d$ robust across five orders of magnitude. We recommend reporting both raw CAR and CAR$_\text{f}$; the gap quantifies action-collapse contamination. This contamination is not architecture-specific: Pi0 Full-FT (${\sim}$2.5B trainable), Pi0 LoRA (${\sim}$33.5M), and Pi0+A1 (${\sim}$658K) all exhibit $>$99\% fake safety despite a 4000$\times$ parameter range. It is a necessary consequence of evaluating collision avoidance without verifying that the robot moved.

More broadly, our results suggest that safety in VLA manipulation may need to be addressed at the planning or constraint-satisfaction layer rather than through auxiliary losses on action prediction. When the action space is EE-centric, any loss defined over EE geometry inherits the same blind spot. Joint-space or full-body constraints, whether learned or imposed at planning time, offer a more natural interface for whole-arm collision avoidance.

\paragraph{Limitations.}
Our evaluation covers one benchmark suite (SafeLIBERO object\_pudding, four tasks) and two model families (Pi0, Pi0.5). Preliminary spatial\_bowl results show different collision patterns, suggesting workspace geometry strongly affects collision modes; generalization to other tasks and architectures (OpenVLA, RT-2) is untested. Oracle experiments were conducted on a largely frozen model (Pi0+A1, TSR ${\approx}$ 0\%); testing on an actively moving policy would better isolate the EE-only structural limitation. The $\lambda = 0.1$ condition reports 2-seed results at the 20K checkpoint; \placeholder{30K three-seed confirmation pending}. Full-body SDF, the natural remedy for the structural mismatch identified here, remains untested. Fine-tuning degrades TSR by ${\sim}$21\textendash23\,pp; strategies that preserve task capability alongside avoidance remain an open problem.

\paragraph{Future work.}
Three directions follow from this analysis: (1) full-body SDF supervision along the kinematic chain, covering intermediate arm links and joints; (2) cross-architecture evaluation on OpenVLA, RT-2, and other VLA families to determine the prevalence of fake safety and metric contamination; (3) planning-layer safety constraints that operate in joint space rather than EE space, decoupling collision avoidance from the action prediction loss.

\paragraph{Broader impact.}
Contaminated CAR scores allow models with no genuine avoidance to pass safety benchmarks, creating false deployment readiness. CAR$_\text{f}$ exposes this gap with minimal changes to existing evaluation protocols.
